\section*{Chapitre \ref{ch:hf:auctions-and-the-close} --- Enchères et clôture}

\begin{solution}{exo:hf:auctions-and-the-close:1}
L'offre atteint 30 actions à 102 (10 à 101 plus 20 à 102)~: l'enchère clôture à 102.
\end{solution}

\begin{solution}{exo:hf:auctions-and-the-close:2}
$9-0.15\times14=6.9$ cents par action~; 8\,280~\$ sur 120\,000 actions.
\end{solution}

\begin{solution}{exo:hf:auctions-and-the-close:3}
Avant 15 h 55~; chaque seconde à partir de 15 h 55.
\end{solution}

\begin{solution}{exo:hf:auctions-and-the-close:4}
Chaque action vendue en plus fait baisser le cours de clôture pour toutes les actions vendues~; le profit de la dernière action marginale tombe à zéro avant que le déséquilibre ne soit comblé
(120\,000 sur 200\,000 ici), et les autres vendeurs de l'enchère comblent le reste à un prix plus élevé.
\end{solution}

\begin{solution}{exo:hf:auctions-and-the-close:5}
Le profit est le cours de clôture moins la valeur du lendemain~; lorsque la moitié du mouvement de 14 cents subsiste, la valeur du lendemain est supérieure de 7 cents et le fournisseur propose aussi
moins (80\,000 actions), donc son cours de clôture est plus élevé mais sa marge plus faible~: 4 cents par action au lieu de 8.
\end{solution}

\begin{solution}{exo:hf:auctions-and-the-close:6}
Si des ordres de vente tardifs réduisent le déséquilibre, une offre au cours de référence fixerait une clôture trop basse pour que le fournisseur réalise un profit~; un ordre à cours limité fixé 4 cents au-dessus
écarte le fournisseur de ces enchères, et coûte peu lorsque le déséquilibre se maintient~: 6.0 cents par action contre 5.5.
\end{solution}

\begin{solution}{exo:hf:auctions-and-the-close:7}
Le déséquilibre fait bouger le cours de clôture de 10 cents au lieu de 14~; la meilleure offre du fournisseur porte sur 100\,000 actions au cours de référence, pour 5.5 cents par action
(contre 6.9)~: un carnet plus profond laisse moins à gagner.
\end{solution}

\begin{solution}{exo:hf:auctions-and-the-close:8}
La moyenne masque les journées où le déséquilibre porte une information~: la plus grande part du mouvement subsiste alors, et une compensation de pleine taille perd de l'argent. Le fournisseur doit dimensionner ses ordres en fonction
du profit espéré, lequel diminue à mesure que la part qui subsiste augmente, et réduire sa taille lorsque le déséquilibre coïncide avec une annonce.
\end{solution}

\begin{solution}{pb:hf:auctions-and-the-close:1}
\begin{enumerate}
\item Les fonds indiciels, les gérants indexés sur un indice et les produits dérivés se règlent au cours de clôture.
\item 7.5~\% du volume quotidien en 2018 (3.1~\% en 2010)~; les cours de clôture se situent généralement au cours acheteur ou au cours vendeur de pré-clôture~; impact plus faible que dans le trading continu~;
les écarts se résorbent rapidement et presque entièrement en moyenne.
\item Publications toutes les 10 secondes à partir de 15 h 50, chaque seconde à partir de 15 h 55~; ordres au marché à la clôture (MOC) avant 15 h 55~; ordres à cours limité à la clôture (LOC) avant 15 h 58~;
modifications avant 15 h 50.
\item Le cours de référence, les actions appariées, le nombre d'actions en déséquilibre et son sens, ainsi que les cours indicatifs proche et éloigné dans les dernières minutes.
\item Voir la \cref{def:hf:auctions-and-the-close:offset}.
\item Voir la \cref{def:hf:auctions-and-the-close:impact}~; 5, 10, 14 et 20 cents pour 25\,000, 100\,000, 200\,000 et 400\,000 actions.
\item $q\,(p-P_0-\pi\Delta(I))$ pour $q$ actions vendues au cours de clôture $p$.
\item Le cours de clôture baisse à chaque action vendue.
\item 120\,000 actions au cours de référence~; le cours de clôture bouge de 9 cents au lieu de 14~; 6.9 cents par action, 8\,280~\$.
\item Elle augmente moins que proportionnellement avec le déséquilibre (4.5, 6.9 et 10 cents pour 100\,000, 200\,000 et 400\,000) et diminue quand $\pi$ augmente (8, 6.9 et 4
cents à 200\,000 pour 0, 15~\% et 50~\%).
\item De 6.9 à 5.5 cents par action~; une prime de 4 cents permet de retrouver 6.0.
\item De 14 cents à 9.
\item Avec 15~\% du mouvement qui subsiste~: 2.25, 2.95, 4.5, 6.9, 8.45 et 10 cents par action pour des déséquilibres de 25\,000 à 400\,000 actions~; 85~\% du
mouvement de clôture se résorbe par hypothèse, et le profit est très sensible à cette hypothèse.
\item Régresser le cours d'ouverture du lendemain sur l'écart de la clôture par rapport au mid de pré-clôture, par titre et par type de journée.
\item Contrats à terme sur indice au règlement-livraison pour la part de marché~; la part idiosyncratique seulement par la taille.
\item Moins d'information, cours de référence périmés~: tailles plus petites, primes plus élevées.
\item Elle la réduit, en ramenant le cours de clôture vers la référence et en absorbant une partie du déséquilibre.
\item Elle rapproche le cours de clôture de la valeur confirmée le lendemain.
\item Fourniture de liquidité à l'enchère de clôture les jours d'annonces.
\item L'immédiateté à la clôture, payée par la part du mouvement qui se résorbe.
\end{enumerate}
\end{solution}

\begin{solution}{iq:hf:auctions-and-the-close:1}
Le déséquilibre par rapport au volume et à la profondeur du titre à la clôture, le cours de référence comparé au carnet continu, toute actualité, les événements d'indice, ainsi que
l'historique des retournements à la clôture du titre, et le budget de risque overnight de la firme.
\iqlookfor{taille relative à la profondeur, actualité, historique des retournements.}
\end{solution}

\begin{solution}{iq:hf:auctions-and-the-close:2}
Pour chaque clôture, l'écart du cours de clôture par rapport au mid de pré-clôture, et le cours du lendemain matin (ou du jour suivant) par rapport à la clôture~; régresser
le mouvement ultérieur sur l'écart, par titre et par période.
\iqlookfor{écart et mouvement ultérieur.}
\end{solution}

\begin{solution}{iq:hf:auctions-and-the-close:3}
Maximiser le volume exécuté, puis minimiser le surplus, puis suivre la pression du marché, puis le cours de référence, avec une règle fixe pour les égalités exactes (le
cours le plus bas, ou celui le plus proche de la référence).
\iqlookfor{hiérarchie des règles.}
\end{solution}

\begin{solution}{iq:hf:auctions-and-the-close:4}
La somme de 800 positions overnight~: une exposition de marché (couverte par des contrats à terme) et 800 écarts idiosyncratiques~; la limiter par une valeur en
risque (VaR) overnight à l'échelle de la firme et par une taille par titre relative au déséquilibre et à la volatilité du titre.
\iqlookfor{agrégation et limites.}
\end{solution}

\begin{solution}{iq:hf:auctions-and-the-close:5}
Les fonds indiciels doivent exécuter le rééquilibrage à la clôture quel que soit le prix, si bien que les déséquilibres sont importants et prévisibles~; les fournisseurs qui les compensent capturent le
retournement.
\iqlookfor{demande contrainte et ses fournisseurs.}
\end{solution}

\begin{solution}{iq:hf:auctions-and-the-close:6}
Le cours de clôture vaut $P_0+\lambda(I-q)$, donc le profit vaut $q(\lambda(I-q)-\pi\lambda I)$, maximal en $q^\ast=I(1-\pi)/2$.
\iqlookfor{condition du premier ordre.}
\end{solution}
